% gtwinxmac.tex -- macros for the master index that gtwill-twinx makes.
% The counterpart of Knuth's twinxmac.tex for CTWILL: the index of a set of
% programs woven by gtwill, set in two columns with gwebmac's index macros. Each
% entry cites the programs by their titles, in the order the files were given:
%     \I{entry}{\twinxref{TITLE}{sections}\twinxref{TITLE2}{sections}...}
% If a file gtwinx-startup.tex can be found, it is read at the end of this file;
% put a localization (\input kotexgweb), a \def\title, or a new \twinxtop there.
\input gwebmac
\catcode`\@=11

% The section numbers belong to other documents, so they link to nothing; a
% definition is underlined, as in each program's own index.
\def\SN#1{#1}
\def\SND#1{$\underline{\hbox{#1}}$}
\let\SNs=\SN \let\SNDs=\SND

\def\title{Index}            % the running head and the title over the index
\def\sechead{}               % there are no sections to show in the running head
\def\grouphead{}
\nosecs                      % no list of section names, no contents
\def\twinxtop{\centerline{\titlefont\title}\vskip 24pt}
\def\twinxsep{}
\def\twinxnext{; }           % between the citations of two programs
\def\twinxref#1#2{\twinxsep{\sc#1}~#2\global\let\twinxsep\twinxnext}
\def\I#1#2{\global\let\twinxsep\empty\II{#1}{#2}}
% A long typewriter string in an entry cannot be broken, and the citations of
% several programs make the lines run differently from any one program's index;
% a hairline overrun is no reason to complain.
\hfuzz=1pt

\openin\z@=gtwinx-startup
\ifeof\z@ \else \closein\z@ \input gtwinx-startup \fi
\catcode`\@=12
